% -------------------------------------------------------
\section{Conclusion}
\label{sec:rs-conclusion}

The Recursive Spawning Protocol makes autonomous drift architecturally impossible — not by prescribing agent behavior, not by enforcing policy conventions, but by establishing immutable parent pointers as a first-class architectural foundation of multi-agent accountability. This is not a minor refinement of existing oversight approaches. It is a structural change to what a delegation chain is: from a mutable reference that can be retroactively rewritten, to a runtime artifact that is immutable from the moment of provisioning and auditable at any future point without trusting any machine actor in the chain.

The protocol achieves this through the mandatory concert of three \bxthree layers — each serving a structurally exclusive role that no other layer can subsume. The Purpose Layer authorizes spawn events and defines the authorized scope refinement policy, establishing the human principal as the exclusive terminus of the delegation chain. The Bounds Engine enforces depth constraints as hard architectural boundaries, preventing unbounded spawn growth that would exceed any human's supervisory capacity. The Fact Layer provides the immutable parent pointer via its write protocol, maintains the tamper-evident forensic ledger, and operates the pre-commit hook that converts the scope refinement constraint from policy documentation into structural enforcement. The three-layer architecture is minimal and sufficient: no proper subset of layers achieves the same guarantees.

The three correctness properties — drift prevention, chain integrity, and termination — are jointly operative and mutually reinforcing. Drift prevention ensures the spawn chain operates only in authorized directions, with each child's operating scope a descendant refinement of its parent's, traceable to the human anchor's intent through the forensic ledger. Chain integrity ensures the chain's topology is stable, its parent pointers unmodifiable after provisioning, and any tampering with the historical record detectable by hash-chain verification. Termination ensures the chain halts within a human-defined depth bound and can sustain no infinite spawn-escalate cycle that circumvents human oversight. Together these properties constitute the RSP's full operational guarantee: a spawn chain is an accountable, bounded, and auditable refinement process that terminates in human judgment.

The limitations analysis demonstrates that RSP's guarantees are not cost-free. The immutable parent pointer prevents legitimate chain reconfiguration; the pre-commit hook and hash-chained forensic ledger impose runtime overhead; the protocol's scope is intra-chain and does not address inter-chain interactions in multi-tenant deployments; and the correctness properties depend on a well-formed Purpose Layer provisioning process that RSP itself does not validate. These are known costs with identified mitigation directions: successor promotion for legitimate reconfiguration, parallel ledger writes for throughput, inter-chain RSP extensions for multi-tenant deployments, and Purpose Layer integrity monitoring for provisioning validation. None of these limitations undermines the core correctness properties under the assumptions stated.

The contribution of RSP to the multi-agent AI accountability landscape is precise and non-overlapping with prior approaches. It is not a monitoring tool added to an existing runtime; it is a structural redefinition of what a delegation chain is. It is not a depth limit with better configuration; it is a proof that depth limits alone cannot prevent drift, and that the conjunction of immutable parent pointers, scope refinement as a structural invariant, and human-exclusive chain termination is both necessary and sufficient for drift prevention. It is not an organizational policy; it is an architectural constraint that holds regardless of the goodness, reliability, or intentions of any machine actor in the chain.

The upstream \bxthree accountability guarantee — that systems built on the \bxthree Framework fail upward into human consciousness, never downward into autonomous chaos — is given its most complete expression in the Recursive Spawning Protocol. When an agent in an RSP-governed system fails, the failure surfaces to a human principal through a complete, immutable, tamper-evident accountability chain. When an RSP-governed system encounters a condition it cannot resolve autonomously, it compulsory-escalates. There is no mechanism by which the system can proceed without human judgment, because the protocol is structured to make human judgment the non-negotiable terminus of every delegation chain.

The Recursive Spawning Protocol is the second named pillar of the \bxthree Framework. It is available for integration into any agentic system that requires the structural guarantee that every agent's authority traces to a human principal — not by policy assertion, not by forensic reconstruction, but by architectural construction from the moment of provisioning.